\chapter{Memory Unsafety}
\label{ch:memory-unsafety}

\epigraph{``A buffer overflow is not a bug. A buffer overflow is a violation of the type system at the machine level.''}{}

\learninggoal{
	\item Define memory unsafety and its root cause in C and C++.
	\item Classify the major families of memory-unsafety vulnerabilities.
	\item Understand stack-based buffer overflows, heap overflows, use-after-free, and format string bugs.
	\item Trace how each vulnerability type leads to a concrete exploitation primitive.
	\item Recognise patterns that indicate unsafe code during code review.
}

\section{What Is Memory Unsafety?}

A programming language is \emph{memory-safe} if every memory access is guaranteed to be within the bounds of the intended object and occurs while that object is still alive. Languages like JavaScript, Python, Java, and Rust (in safe mode) enforce this guarantee---either through runtime checks, garbage collection, or a borrow checker.

C and C++ are \emph{memory-unsafe}. They provide no automatic bounds checking, no lifetime tracking, and no protection against dangling pointers. Once a pointer is computed, the processor will dereference it regardless of whether the target is valid.

\begin{definition}[Memory Safety Violation]
	A \textbf{memory safety violation} occurs when a program accesses memory outside the bounds of the intended object, or accesses memory through a pointer to an object that has already been deallocated. The C standard classifies these as \emph{undefined behaviour}---the compiler may generate any code.
\end{definition}

Memory-safety bugs are not merely quality issues. They are the dominant source of security vulnerabilities in systems software. Microsoft estimates that 70\% of all CVEs in their products are memory-safety issues. Google finds similar numbers in Android and Chrome.

\section{Classification of Memory Unsafety}

We can classify memory unsafety along two axes: \emph{spatial} (accessing memory beyond the object's bounds) and \emph{temporal} (accessing memory after the object's lifetime):

\begin{longtable}{p{3cm}p{5cm}p{5cm}}
	\toprule
	                & \textbf{Spatial} (out-of-bounds) & \textbf{Temporal} (use-after-free)  \\
	\midrule
	\textbf{Stack}  & Stack buffer overflow            & Return of pointer to local variable \\
	\textbf{Heap}   & Heap buffer overflow             & Use-after-free, double-free         \\
	\textbf{Global} & Global buffer overflow           & N/A (global lifetime)               \\
	\bottomrule
\end{longtable}

\section{Stack Buffer Overflows}

The stack buffer overflow is the classic memory-safety vulnerability, documented by Aleph One in 1996~\cite{aleph1996smashing}.

\begin{lstlisting}[language=C, caption=Vulnerable stack buffer]
void vulnerable(char *input) {
    char buf[64];
    strcpy(buf, input);  // No bounds check!
}
\end{lstlisting}

The function allocates 64 bytes on the stack. If \texttt{input} is longer than 63 characters (plus null terminator), \texttt{strcpy} writes past the end of \texttt{buf}, overwriting the saved \texttt{rbp} and the return address.

\begin{principle}[The Exploitation Primitive]
	Overwriting the return address gives the attacker control of the instruction pointer (\texttt{rip}) when the function returns. The attacker redirects execution to their shellcode or a ROP chain.
\end{principle}

The layout just before \texttt{strcpy} writes:

\begin{lstlisting}[caption=Stack layout during overflow]
  High addresses
  +---------------------------+
  | Return address            |  <- overwritten by attacker
  +---------------------------+
  | Saved rbp                 |  <- overwritten
  +---------------------------+
  | buf[63..56]               |  <- top of buffer
  | buf[55..48]               |
  | ...                        |
  | buf[7..0]                 |  <- bottom of buffer (buf[0])
  +---------------------------+  <- rsp
  Low addresses
\end{lstlisting}

\subsection{Functions That Enable Overflows}

Certain C library functions are notorious for lacking bounds checking:

\begin{longtable}{p{3cm}p{4.5cm}p{4.5cm}}
	\toprule
	\textbf{Function}               & \textbf{Problem}                & \textbf{Safe alternative}                           \\
	\midrule
	\texttt{strcpy(dst, src)}       & No length argument              & \texttt{strncpy}, \texttt{strlcpy}, \texttt{memcpy} \\
	\texttt{sprintf(buf, fmt, ...)} & No buffer size                  & \texttt{snprintf}                                   \\
	\texttt{gets(buf)}              & No limit                        & \texttt{fgets}                                      \\
	\texttt{scanf("\%s", buf)}      & No field width                  & \texttt{scanf("\%Ns", buf)}                         \\
	\texttt{memcpy(dst, src, n)}    & Relies on caller for \texttt{n} & Manual bounds check                                 \\
	\bottomrule
\end{longtable}

\section{Heap Overflows}

Heap overflows are analogous to stack overflows but occur in dynamically allocated memory.

\begin{lstlisting}[language=C, caption=Vulnerable heap buffer]
void heap_vulnerable(char *input) {
    char *buf = malloc(64);
    strcpy(buf, input);  // If input > 63 bytes, overflow
}
\end{lstlisting}

Heap exploitation is more complex because the target is not a fixed return address but the heap metadata or adjacent heap objects. Modern heap allocators use complex data structures (bins, fastbins, tcache) that can be corrupted to achieve arbitrary write.

\begin{definition}[Heap Metadata Corruption]
	Heap allocators store metadata (chunk size, flags, free-list pointers) in-band with user data. Overwriting this metadata can cause the allocator to write attacker-controlled values to attacker-controlled addresses---a \emph{write-what-where} primitive.
\end{definition}

A typical heap chunk layout (in \texttt{ptmalloc}):

\begin{lstlisting}[caption=Heap chunk structure (simplified)]
  Low address
  +---------------------------+
  | prev_size (8 bytes)       |  <- size of previous chunk
  +---------------------------+
  | size (8 bytes)            |  <- current chunk size + flags
  +---------------------------+
  | user data ...             |  <- what malloc() returns
  | ...                        |
  +---------------------------+
  High address
\end{lstlisting}

Overflowing from chunk A into chunk B can corrupt B's size field or, if B is freed, its free-list pointers. When B is later allocated, the allocator dereferences the corrupted pointers, giving the attacker a controlled write.

\section{Use-After-Free}

A \textbf{use-after-free} (UAF) bug occurs when memory is accessed through a pointer after it has been freed.

\begin{lstlisting}[language=C, caption=Use-after-free example]
struct object *obj = malloc(sizeof(*obj));
// ... use obj ...
free(obj);
// ... obj is now a dangling pointer ...
obj->field = 42;  // Use-after-free! Writing to freed memory
\end{lstlisting}

After \texttt{free}, the memory at \texttt{obj} may be reused by a subsequent allocation. If an attacker can control the data in that new allocation (e.g., by allocating a buffer of the same size with attacker-controlled content), they can manipulate the original object's fields.

\begin{definition}[Dangling Pointer]
	A \textbf{dangling pointer} is a pointer that references memory that has been freed. Dereferencing a dangling pointer is undefined behaviour that typically results in a use-after-free vulnerability.
\end{definition}

UAF exploitation often involves:

\begin{enumerate}
	\item Allocating an object of type \texttt{A}.
	\item Freeing it while retaining a pointer.
	\item Allocating a controlled buffer of the same size (which gets the freed \texttt{A}'s memory).
	\item Accessing the original pointer, which now reads/writes attacker-controlled data.
\end{enumerate}

\subsection{Double Free}

A related temporal bug is the \textbf{double free}: freeing the same pointer twice. This corrupts the allocator's free-list data structures.

\begin{lstlisting}[language=C, caption=Double free]
void *p = malloc(64);
free(p);
free(p);  // Double free!
\end{lstlisting}

After the first \texttt{free}, the allocator adds \texttt{p} to the free list. After the second \texttt{free}, \texttt{p} is added again. Subsequent allocations may return the same memory twice, bypassing the allocator's integrity guarantees.

\section{Format String Vulnerabilities}

Format string bugs are an entirely different class of memory unsafety.

\begin{lstlisting}[language=C, caption=Format string vulnerability]
void vulnerable(char *user_input) {
    printf(user_input);  // Attacker controls the format string!
}
\end{lstlisting}

The correct call is \texttt{printf("\%s", user\_input)}. When the attacker supplies a format string like \texttt{"\%x \%x \%x \%x \%x"}, \texttt{printf} reads arguments from the stack that were never passed, leaking memory contents.

More dangerously, the \texttt{\%n} specifier writes the number of bytes printed so far to a pointer argument:

\begin{lstlisting}[caption=Using \%n to write memory]
  // If the attacker can supply "%n" to printf,
  // printf will write 4 bytes to the address on the stack
  // where it expects the next argument.
  printf("%x%x%x%n");  // writes to stack, corrupting memory
\end{lstlisting}

This gives the attacker an arbitrary-write primitive by carefully controlling the number of printed bytes and the target address (planted on the stack by preceding format arguments).

\begin{longtable}{p{2.5cm}p{4cm}p{5cm}}
	\toprule
	\textbf{Specifier} & \textbf{Effect}                 & \textbf{Security relevance}      \\
	\midrule
	\texttt{\%x}       & Read unsigned hex from stack    & Stack memory leak                \\
	\texttt{\%s}       & Read string from stack pointer  & Read memory at arbitrary address \\
	\texttt{\%p}       & Read pointer from stack         & Leak addresses (bypass ASLR)     \\
	\texttt{\%n}       & Write count to pointer on stack & Arbitrary memory write           \\
	\texttt{\%hn}      & Write 2 bytes (short)           & Fine-grained memory write        \\
	\texttt{\%hhn}     & Write 1 byte (char)             & Byte-level memory write          \\
	\bottomrule
\end{longtable}

\section{Integer Overflows Leading to Memory Unsafety}

Integer overflows are safety bugs that often mask memory-safety vulnerabilities:

\begin{lstlisting}[language=C, caption=Integer overflow enabling overflow]
void vulnerable(size_t size) {
    char *buf = malloc(size + 1);  // If size = SIZE_MAX, size+1 = 0!
    read(fd, buf, size);           // Overflow into adjacent heap memory
}
\end{lstlisting}

If \texttt{size} is \texttt{SIZE\_MAX}, then \texttt{size + 1} wraps to 0. \texttt{malloc(0)} returns a valid (but tiny) buffer. The subsequent \texttt{read} writes \texttt{SIZE\_MAX} bytes into that small buffer.

\section{The Eternal War}

Szekeres et al.~\cite{szekeres2013eternal} argue that memory corruption is an \emph{undecidable} problem in general: no static analysis can prove the absence of all memory-safety bugs in arbitrary C/C++ programs. The best we can do is:

\begin{longtable}{p{4cm}p{4.5cm}p{4.5cm}}
	\toprule
	\textbf{Approach}     & \textbf{Mechanism}                          & \textbf{Limitation}                               \\
	\midrule
	Static analysis       & Symbolic execution, abstract interpretation & False positives, incomplete coverage              \\
	Dynamic analysis      & AddressSanitizer, valgrind                  & Runtime overhead, no coverage of unexecuted paths \\
	Memory-safe languages & Type systems, borrow checker                & Runtime overhead (GC), ecosystem friction         \\
	Hardware protection   & NX, ASLR, MPK                               & Doesn't prevent the bug, only limits exploitation \\
	\bottomrule
\end{longtable}

\section{Exercises}

\begin{exercise}
	Identify the memory-safety bug in this code. What is the vulnerability class? How would you fix it?
	\begin{lstlisting}[language=C]
void copy_name(char *name) {
    char buf[16];
    strcpy(buf, name);
    printf("Name: %s\n", buf);
}
	\end{lstlisting}
\end{exercise}

\begin{exercise}
	What is the output of this program? What security implication does it have?
	\begin{lstlisting}[language=C]
int main() {
    char *s = malloc(8);
    free(s);
    printf("%p\n", s);
}
	\end{lstlisting}
\end{exercise}

\begin{exercise}
	Given \texttt{printf(user\_input)}, how would an attacker leak a secret value at address \texttt{0x7fffffffe000} using format string specifiers?
\end{exercise}

\begin{exercise}
	Research CVE-2021-26708 (a Linux kernel use-after-free). Describe the bug, the exploit technique, and the patch. What allocation type was reused?
\end{exercise}
